% =====================================================================
%  Shadrach
%  A Cut-Primitive Domain-Specific Language for Financial Instruments
%  by Individuation
%
%  Self-contained. Every mathematical result the language depends on
%  is derived or stated in full within this document; no companion
%  manuscript is assumed. A reader who reads only this paper has
%  everything needed to understand and implement the language.
% =====================================================================
\documentclass[11pt,a4paper]{article}

\usepackage[utf8]{inputenc}
\usepackage[T1]{fontenc}
\usepackage{lmodern}
\usepackage[a4paper,margin=2.4cm]{geometry}
\usepackage{amsmath,amssymb,amsthm,mathtools}
\usepackage{bm}
\usepackage{mathrsfs}
\usepackage{booktabs}
\usepackage{array}
\usepackage{enumitem}
\usepackage{microtype}
\usepackage{xcolor}
\usepackage{listings}
\usepackage{algorithm}
\usepackage{algpseudocode}
\usepackage[round,sort&compress]{natbib}
\usepackage[colorlinks=true,linkcolor=blue!45!black,
            citecolor=blue!45!black,urlcolor=blue!45!black]{hyperref}
\usepackage{cleveref}

% ---- Theorem-like environments ----
\newtheoremstyle{thmstyle}{\topsep}{\topsep}{\itshape}{}{\bfseries}{.}{.5em}{}
\theoremstyle{thmstyle}
\newtheorem{theorem}{Theorem}[section]
\newtheorem{lemma}[theorem]{Lemma}
\newtheorem{proposition}[theorem]{Proposition}
\newtheorem{corollary}[theorem]{Corollary}
\theoremstyle{definition}
\newtheorem{definition}[theorem]{Definition}
\newtheorem{rulename}[theorem]{Rule}
\newtheorem{principle}[theorem]{Principle}
\newtheorem{construction}[theorem]{Construction}
\newtheorem{example}[theorem]{Example}
\theoremstyle{remark}
\newtheorem{remark}[theorem]{Remark}
\newtheorem{notation}[theorem]{Notation}

% ---- shadrach listing style ----
\definecolor{shkw}{RGB}{30,64,175}
\definecolor{shcom}{RGB}{22,101,52}
\definecolor{shnum}{RGB}{180,83,9}
\definecolor{shstr}{RGB}{159,18,57}
\definecolor{shbg}{RGB}{250,250,250}
\lstdefinelanguage{shadrach}{
  morekeywords={floor,cut,bond,close,track,until,yield,when,emit,observe,
    where,in,as,let,medium,converge,diverge,with,by,reps,
    representation,switch,import,module,export,assert,from,to,if,else,
    position,portfolio,rebalance},
  morekeywords=[2]{Position,Bond,Portfolio,Path,Scalar,price,qty,corr,
    weight,vacancy,floorval,residual},
  sensitive=true,
  morecomment=[l]{--},
  morestring=[b]",
  keywordstyle=\color{shkw}\bfseries,
  keywordstyle=[2]\color{shnum},
  commentstyle=\color{shcom}\itshape,
  stringstyle=\color{shstr},
  basicstyle=\ttfamily\small,
  backgroundcolor=\color{shbg},
  frame=single,framesep=4pt,xleftmargin=6pt,xrightmargin=4pt,
  breaklines=true,showstringspaces=false,columns=fullflexible,tabsize=2
}
\lstset{language=shadrach}

\crefname{rulename}{Rule}{Rules}
\Crefname{rulename}{Rule}{Rules}

% ---- shorthand ----
\newcommand{\R}{\mathbb{R}}
\newcommand{\N}{\mathbb{N}}
\newcommand{\Z}{\mathbb{Z}}
\newcommand{\Gph}{\Gamma}
\newcommand{\V}{V}
\newcommand{\E}{E}
\newcommand{\wt}{w}
\newcommand{\flo}{\beta}
\newcommand{\res}{\rho}
\newcommand{\cut}{\partial}
\newcommand{\co}{\complement}
\newcommand{\medium}{\mathsf{m}}
\newcommand{\evalto}{\Downarrow}
\newcommand{\dd}{\,\mathrm{d}}
\newcommand{\bnf}{\;::=\;}
\newcommand{\alt}{\;\mid\;}
\newcommand{\sh}[1]{\texttt{#1}}
\newcommand{\abs}[1]{\lvert#1\rvert}
\newcommand{\norm}[1]{\lVert#1\rVert}
\newcommand{\Lap}{\mathcal{L}}
\newcommand{\simplex}{\Delta_m}
\newcommand{\Rec}{\mathcal{R}}
\newcommand{\eps}{\varepsilon}
\DeclareMathOperator*{\argmax}{arg\,max}

\title{\textbf{Shadrach:\\[3pt]
A Cut-Primitive Domain-Specific Language for\\
Financial Instruments by Individuation}}

\author{%
Kundai Farai Sachikonye\\
\small Technical University of Munich\\
\small AIMe Registry for Artificial Intelligence\\
\small \texttt{kundai.sachikonye@wzw.tum.de}}

\date{\today}

\begin{document}
\maketitle

% =====================================================================
\begin{abstract}
\noindent
We specify \emph{Shadrach} (source extension \sh{.shad}), a
domain-specific language for financial computation in which the sole
computational primitive is the \emph{cut}: the act of individuating a
part from the rest of a bounded whole at a strictly positive,
irreducible cost --- the \emph{floor} $\flo>0$. The language rests on a
single derivation, given here in full: a price is never a point but a
bounded \emph{cell}, because comparing an instrument against the
entirety of the market that prices it is a comparison that can never be
completed, and the bid--ask spread is the irreducible residue of that
incompleteness, not a friction to be modelled away. From this one fact
we build a single semantic identity, carried through the whole
language: opening a position, correlating two positions, closing a
portfolio to full allocation, and rebalancing a portfolio toward
equilibrium are not four distinct operations but the same cut at four
different arities. A length-one cut individuates a \emph{position}; a
cut between two positions is a \emph{bond}, existing only if it lowers
total un-allocated capital and carrying a Kirchhoff conductance fixed by
the pair's return covariance; a portfolio is a \emph{close}: a set of
bonds driving every position's un-allocated capital (\emph{vacancy}) to
zero, i.e.\ Kirchhoff's current law satisfied at every node; and a
\emph{rebalance} is a \emph{track}: a residue-chained sequence of cuts
converging, by the Banach fixed-point theorem, to the unique
equilibrium allocation. Shadrach therefore has one verb and one
composition rule, from which position-opening, correlation, portfolio
construction, and equilibrium rebalancing all follow. We give: (i)~the
semantic model, a finite weighted \emph{contact graph} whose distinguished
vertex is the market medium and on which every Shadrach value is a cut
of accountable residue, never a bare number; (ii)~a full derivation, from
two premises alone, of the positive floor and of individuation by
negation, so that no theorem below is posited; (iii)~a complete lexical
and context-free grammar in EBNF; (iv)~a static \emph{accountability}
type system in which every value carries the floor at which it was cut,
so a value claiming zero residue --- a frictionless, perfectly liquid
market --- is rejected at compile time; (v)~a small-step operational
semantics in which evaluation \emph{is} measurement, with a monotone
committed-cut count as an intrinsic, non-decreasing clock; (vi)~a
standard library binding each verb to a named theorem proved in this
paper --- \sh{cut} to the Price-Cell theorem, \sh{bond} to a Kirchhoff
conductance criterion, \sh{close} to the Kirchhoff Portfolio Formula,
and \sh{track} to Banach-contraction rebalancing --- each restated and
re-proved here from first principles, with the numerical values every
implementation must reproduce; and (vii)~a two-target compilation
model: a Rust reference implementation realising cuts as exact
minimum-cut and pseudoinverse computations, and a TypeScript/WebGL\,2
interactive runtime realising cuts as GPU render passes, proved
observationally equivalent up to the floor. The specification is
implementable as written.
\end{abstract}

\noindent\textbf{Keywords:} domain-specific language; quantitative
finance; individuation; cut primitive; contact graph; Kirchhoff's laws;
Banach fixed point; accountability type system; operational semantics;
portfolio construction; ETF rebalancing.

\tableofcontents

% =====================================================================
\part{Foundations: why a price is a cut}
% =====================================================================

% =====================================================================
\section{Introduction}
\label{sec:intro}
% =====================================================================

\subsection{The question}

What does it mean to price a financial instrument, and what is a
portfolio? The working assumption of most quantitative software is that
a price is, in the limit of perfect information, a single real number:
frictions --- the bid--ask spread, market impact, information asymmetry
--- are treated as measurement error superimposed on an underlying exact
value, and a portfolio is an accounting artefact, a vector of such
numbers weighted by holdings. This paper takes the opposite position and
proves it. A price is not a number with error attached; it is,
irreducibly, a \emph{bounded region} --- and the width of that region
does not shrink to zero as information improves, because the comparison
that would shrink it to zero can never be completed. A portfolio is not
a vector of independent numbers; it is a graph in which capital
conservation (Kirchhoff's current law) and no-arbitrage (Kirchhoff's
voltage law) are structural constraints, and its equilibrium allocation
is not chosen but computed, as the unique fixed point of a provably
contracting map.

This is not a modelling choice; it is derived, in \Cref{sec:ground}
below, from two facts that require no financial assumption at all: that
an instrument is a proper part of the market that prices it, and that
the market is never exhausted by any finite act of price discovery. The
present paper makes this derivation literal by defining a language,
\emph{Shadrach} (\sh{.shad}), whose every statement is one of the
resulting theorems, executed. A financial model becomes a script and an
equilibrium computation becomes a run.

\subsection{One primitive: the cut}

Shadrach has a single computational primitive, the \emph{cut}, for
exactly the reason given above: individuating any part of a bounded
whole --- separating one instrument's value from the undifferentiated
mass of the market that prices it --- costs a strictly positive,
irreducible minimum, the floor $\flo>0$. Every Shadrach operation is a
cut:

\begin{itemize}[leftmargin=1.4em,itemsep=2pt]
\item A \emph{position} is a cut of arity one: individuating a holding
in instrument $i$ from the market medium (\Cref{sec:stdlib-position}).
\item A \emph{bond} is a cut of arity two: the shared boundary between
two positions, existing iff a Kirchhoff conductance derived from their
return covariance is positive (\Cref{sec:stdlib-bond}).
\item A \emph{portfolio} is a closure of cuts: cuts continued until
every position's un-allocated capital (\emph{vacancy}) is driven to
zero --- Kirchhoff's current law satisfied everywhere
(\Cref{sec:stdlib-portfolio}).
\item A \emph{rebalance} is a residue-chained sequence of cuts: the
residue of each cut is the input to the next, converging geometrically
to the unique equilibrium allocation (\Cref{sec:stdlib-track}).
\end{itemize}

Opening a position is not a different verb from rebalancing a
portfolio; it is a track of length one. This is forced by the theory,
not chosen for economy of notation: in the framework developed below, a
trade and a price observation are the same cut, and position-opening,
correlation, portfolio construction, and rebalancing are that cut at
arities one, two, closure, and chain. Shadrach therefore has one verb
(\sh{cut}) and one composition rule (residue chaining); every other
surface form is sugar over these two.

\subsection{Why this document is self-contained}

Everything Shadrach needs --- the positive floor, individuation by
negation, the contact-graph model of a market, the Kirchhoff conductance
criterion for bonding, the closed-form portfolio formula, and the
Banach-contraction argument for rebalancing --- is derived or restated
in full below, from two premises (\Cref{sec:ground}) and standard,
externally published graph theory, linear algebra, and fixed-point
theory. No result is imported by reference to a companion manuscript. A
reader who accepts \Cref{def:whole,def:noncompletable} (an instrument is
a proper part of a market that is never exhausted by finite price
discovery) can verify every theorem that follows without consulting any
other document.

\subsection{Organisation}

\Cref{sec:ground} derives the positive floor from two premises about
markets and instruments alone. \Cref{sec:model} fixes the semantic
model: the market contact graph, the cut, the floor, the residue.
\Cref{sec:negation} proves individuation by negation for instruments.
\Cref{sec:lexical,sec:grammar} give the lexical and context-free
grammar. \Cref{sec:types} defines the accountability type system and
proves its soundness. \Cref{sec:opsem} gives the small-step operational
semantics and proves cut monotonicity. \Cref{sec:stdlib} specifies the
standard library, proving from scratch the Price-Cell theorem, the
Kirchhoff conductance bonding criterion, the Kirchhoff Portfolio
Formula, and Banach-contraction rebalancing, each with the exact
numerical predictions an implementation must reproduce.
\Cref{sec:ir} defines the core intermediate representation and the
two-target compilation model, and proves target equivalence.
\Cref{sec:examples} gives worked programs. \Cref{sec:impl} states
implementation requirements. \Cref{sec:validation} reports the
numerical validation suite. \Cref{sec:conclusion} concludes.

% =====================================================================
\section{The ground: why the floor is derived, not assumed}
\label{sec:ground}
% =====================================================================

We derive the central constant of the theory, the floor $\flo>0$, from
two premises that mention no financial mechanism: an instrument is a
proper part of a whole, and the whole is never exhausted by any finite
act of comparison. Everything downstream --- the impossibility of a
frictionless price, the necessity of a bid--ask spread, the boundedness
of every portfolio allocation --- is a consequence, not a postulate.

\begin{definition}[The market; an instrument]
\label{def:whole}
The \emph{market} $\medium$ is the totality of information, capital, and
counterparties that could bear on the value of a financial instrument.
An \emph{instrument} is a part $A$ of $\medium$ that is \emph{proper}:
$A\neq\medium$ and $A\neq\varnothing$. An instrument is, by definition,
not the market that prices it.
\end{definition}

\begin{definition}[Non-completable market]
\label{def:noncompletable}
The market $\medium$ is \emph{non-completable} if no finite stage of
price discovery exhausts it: for every finite set $F$ of facts,
counterparties, and transactions already brought to bear on an
instrument's valuation, there exists a further fact, counterparty, or
potential transaction not in $F$. This is an order condition --- no
last stage of price discovery --- not a claim about the cardinality of
the market.
\end{definition}

\begin{remark}[Why markets are non-completable]
\label{rem:why-noncompletable}
For any finite state of information about an instrument --- its filings,
its order book, its analyst coverage, its correlated instruments'
prices --- there remains a further counterparty not yet transacted with,
a further piece of private information not yet disclosed, a further
future cash flow not yet realised. This is not a claim about how much
information exists, only that no finite act of price discovery brings
all of it to bear. That is exactly \Cref{def:noncompletable}.
\end{remark}

\begin{definition}[Identification; identification residual]
\label{def:identification}
To \emph{identify} (price) an instrument $A$ is to compare it against
the rest of the market, $\medium\setminus A$. The identification is
\emph{complete} if every fact of the market has entered the comparison.
The \emph{identification residual} $\res(A)\ge0$ is the content of the
comparison that remains unfinished; $\res(A)=0$ iff the identification
is complete.
\end{definition}

\begin{theorem}[Floor from Infinitude]
\label{thm:floor-infinitude}
Let $\medium$ be non-completable and let $A$ be an instrument (a proper
part, \Cref{def:whole}) identified by comparison against
$\medium\setminus A$. Then the identification of $A$ cannot be
completed, and
\[
\boxed{\ \res(A)\ >\ 0.\ }
\]
\end{theorem}

\begin{proof}
Suppose $\res(A)=0$: the identification is complete, so by
\Cref{def:identification} every fact of $\medium\setminus A$ has been
brought to bear. Since $A$ is fixed, bringing the whole of
$\medium\setminus A$ to bear exhausts $\medium$ up to the fixed part
$A$ --- a finite addition to a finite comparison. This produces a finite
stage that exhausts $\medium$, contradicting non-completability
(\Cref{def:noncompletable}). Hence no finite comparison completes the
identification of $A$, and $\res(A)>0$.
\end{proof}

\begin{corollary}[The floor is derived]
\label{cor:floor-derived}
There is a constant $\flo:=\inf_A\res(A)>0$, the \emph{floor}, bounding
below the identification residual of every instrument. It is not an
independent assumption but a consequence of: (i) an instrument is a
proper part of the market, and (ii) the market is non-completable.
\end{corollary}

\begin{proof}
By \Cref{thm:floor-infinitude}, $\res(A)>0$ for every instrument $A$. If
$\inf_A\res(A)=0$ there would be instruments with arbitrarily
near-complete identifications; in the limit such an identification would
finish the comparison and exhaust $\medium$, the same contradiction as in
\Cref{thm:floor-infinitude}. Hence $\flo:=\inf_A\res(A)>0$.
\end{proof}

\begin{remark}[The bid--ask spread is the floor, not a friction]
\label{rem:spread-is-floor}
\Cref{cor:floor-derived} is the precise sense in which a bid--ask spread
is not a market imperfection to be modelled away by better information
or lower latency. It is the trace the non-completable market leaves on
every finite instrument: the residue of a comparison that cannot be
finished. A ``perfect'' market --- zero spread, a price resolved to a
point --- would require completing a comparison against a market that,
by hypothesis, never completes. \Cref{sec:stdlib-position} makes this
literal: every \sh{Position} value in Shadrach carries a floor
annotation $\flo>0$ that no computation can drive to zero.
\end{remark}

% =====================================================================
\section{The semantic model}
\label{sec:model}
% =====================================================================

Shadrach is defined over a single mathematical object, the market
contact graph, and a single operation on it, the cut.

\subsection{The market contact graph}

\begin{definition}[Contact graph]
\label{def:contact-graph}
A \emph{contact graph} is a finite weighted graph $\Gph=(\V,\E,\wt)$
with a distinguished \emph{medium} vertex $\medium\in\V$ adjacent to
every other vertex, and $\wt\colon\E\to\R_{>0}$ a strictly positive
weight. Vertices $v\neq\medium$ are \emph{positions}: individuated
holdings in a financial instrument. An edge between two positions is a
\emph{contact}; its weight is the cost of the distinction maintained
between them. $\Gph$ is the runtime state of a Shadrach program: every
value the program holds is a structure on $\Gph$.
\end{definition}

\begin{definition}[Floor]
\label{def:floor}
The \emph{floor} of a contact graph is a constant $\flo>0$ with
$\wt(e)\ge\flo$ for every edge $e\in\E$ --- the constant derived in
\Cref{cor:floor-derived}. No position is individuated from the market,
and no two positions are bonded, at zero cost. The floor is a
program-level parameter (set by \sh{floor}, \Cref{sec:grammar}); its
positivity is enforced by the type system (\Cref{sec:types}).
\end{definition}

\begin{definition}[Cut; residue]
\label{def:cut}
A \emph{cut} of a position set $S$ ($\varnothing\neq S\subsetneq\V$,
$\medium\notin S$) is the edge boundary
$\cut(S)=\{e\in\E:\abs{e\cap S}=1\}$, of weight
$\wt(\cut(S))=\sum_{e\in\cut(S)}\wt(e)\ge\flo$. The cut individuates $S$
from the rest of the graph. A cut \emph{event} commits one such
boundary, fixing its edges as realised separations and advancing the
program's cut count by one. The \emph{residue} of a cut event is the
weight $\res=\wt(\cut(S))\ge\flo$ it deposits: the value the cut
produces, and the boundary condition of the next cut in a chain
(\Cref{sec:stdlib-track}).
\end{definition}

\begin{principle}[Accountability]
\label{prin:accountability}
Every Shadrach value is a cut, and every cut carries its residue. There
are no bare numbers in Shadrach: a price, a weight, a risk figure is
always paired with the floor at which it was individuated. A value
claiming residue below the floor, or exactly zero --- a frictionless
market, a costless bond, a fully arbitraged portfolio with no residual
uncertainty --- is not a value; it is the forbidden sharp cut, and the
type system rejects it (\Cref{thm:no-zero-value}).
\end{principle}

\begin{remark}[Why a graph and not a number line]
\label{rem:why-graph}
Shadrach computes on a contact graph rather than on real numbers because
the quantities it manipulates --- an instrument's price cell, a pair's
bonding conductance, a portfolio's equilibrium allocation --- are
minimum cuts and graph Laplacian quantities, not fitted scalars. A
position's price is the weight of its cut against the market medium; a
bond's conductance is fixed by the covariance structure of the pair it
joins; a portfolio's equilibrium is the solution of a Kirchhoff system
on the whole graph. Representing these as graph cuts, rather than as
independent numbers with error bars attached, is what lets the same
primitive serve position-opening, bonding, portfolio-closure, and
rebalancing (\Cref{rem:one-verb}).
\end{remark}

\subsection{Measurement, not simulation}

\begin{principle}[Evaluation is measurement]
\label{prin:measure}
Executing a Shadrach cut performs the individuation; it does not
simulate a description of it. The result of \sh{cut AAPL qty 100} is
the measured price cell of a 100-share position in AAPL against the
current market medium, carried as an accountable value, not a fitted
model of it. Operationally (\Cref{sec:opsem}) a cut evaluation mutates
the contact graph --- it commits boundary and advances the cut count ---
rather than returning a pure function of its inputs. Re-evaluating a
cut is a new cut at a higher count, never a cached recomputation
(\Cref{thm:monotone}).
\end{principle}

% =====================================================================
\section{Individuation by negation}
\label{sec:negation}
% =====================================================================

The floor says every individuation costs something; it does not yet say
\emph{how} a position is determined. We show a position can be fixed
only by its complement --- by the market it is not part of --- and that
any attempt to fix it by naming a distinguished, privileged instrument
in advance regresses without end.

\begin{definition}[Selector; no privileged instrument]
\label{def:selector}
A \emph{selector} for a position set $U\subseteq\V$ is a rule that fixes
$U$ by naming a distinguished vertex of $U$ in advance of any market
data. A determination of $U$ is \emph{selector-free} if it names no
member of $U$. The market contact graph has \emph{no privileged
instrument} if the data available to a determining rule are invariant
under the weighted automorphism group $\mathrm{Aut}(\Gph)$: no
instrument is singled out by market structure alone, prior to trading.
\end{definition}

\begin{theorem}[Individuation by negation]
\label{thm:negation}
Let $\Gph$ have no privileged instrument. Then any selector-free
determination of a position set $U$ fixes $U$ as the complement of its
negation set, $U=\V\setminus\co U$ with $\co U:=\V\setminus U$ specified
first; conversely every admissible $\co U$ determines a unique $U$; and
no selector-based rule fixes $U$ without invoking a further selector.
\end{theorem}

\begin{proof}
\emph{Necessity.} To fix $U$ is to decide, for every vertex of $\Gph$,
whether it belongs to $U$ or its complement. A \emph{positive} rule ---
one asserting membership of a distinguished vertex of $U$ --- must name
that vertex; since the data are $\mathrm{Aut}(\Gph)$-invariant
(\Cref{def:selector}), no vertex is named by the data alone, so the rule
itself must supply a selector, which must in turn be fixed --- either
given in advance (excluded) or fixed by a prior rule facing the same
dichotomy. The regress terminates only at a rule naming no member of
$U$: specify the complement $\co U=\V\setminus U$ and take $U$ as what
remains. This names only ``the rest of the market,'' a datum already
present once $\V$ is given. \emph{Sufficiency and uniqueness.}
Complementation $U\mapsto\V\setminus U$ is an involution on the finite
set $\V$, so $\co U$ determines $U$ uniquely and conversely.
\end{proof}

\begin{corollary}[A position is fixed by what it is not exposed to]
\label{cor:position-negation}
A single position $v$ is determined selector-free by the partition of
the rest of the market into its correlated instruments $N(v)$ (what it
bonds to) and its uncorrelated complement $\V\setminus N[v]$ (what it
neither is nor is exposed to). No intrinsic label distinguishes $v$
in advance of this negation structure.
\end{corollary}

\begin{remark}[Why this matters for pricing]
\label{rem:negation-pricing}
\Cref{thm:negation} is the formal reason a Shadrach program's mandatory
\sh{not\{\}} clause (\Cref{sec:grammar}) is not stylistic. An instrument
priced without reference to what it is \emph{not} exposed to --- its
excluded risk factors, its uncorrelated universe --- has not been
individuated at all in a market with no privileged instrument; the
market alone does not single out any instrument as the answer.
\end{remark}

% =====================================================================
\part{Syntax}
% =====================================================================

% =====================================================================
\section{Lexical structure}
\label{sec:lexical}
% =====================================================================

A Shadrach source file is UTF-8 text with extension \sh{.shad}. Lexing
produces a token stream from the following classes.

\begin{definition}[Token classes]
\label{def:tokens}
\leavevmode
\begin{itemize}[leftmargin=1.4em,itemsep=1pt]
\item \textbf{Comments.} \sh{--} to end of line; discarded.
\item \textbf{Keywords.} \sh{floor}, \sh{cut}, \sh{bond}, \sh{close},
\sh{track}, \sh{until}, \sh{yield}, \sh{when}, \sh{do}, \sh{emit},
\sh{observe}, \sh{in}, \sh{as}, \sh{let}, \sh{medium}, \sh{converge},
\sh{diverge}, \sh{with}, \sh{by}, \sh{import}, \sh{module}, \sh{export},
\sh{assert}, \sh{position}, \sh{portfolio}, \sh{rebalance}.
\item \textbf{Identifiers.} \sh{[A-Za-z\_][A-Za-z0-9\_]*}, not a
keyword. Instrument tickers (e.g.\ \sh{AAPL}) are identifiers.
\item \textbf{Numbers.} decimal integers and floats with optional
exponent: \sh{[0-9]+(.[0-9]+)?([eE][+-]?[0-9]+)?}.
\item \textbf{Strings.} double-quoted, with the usual escapes.
\item \textbf{Operators and punctuation.} \sh{:=} (bind), \sh{\~{}}
(cut between / bond), \sh{( ) [ ] \{ \} , : . > < >= <= == ->} and the
suffix forms used by literals.
\end{itemize}
Whitespace is insignificant except as a token separator; Shadrach is not
layout-sensitive.
\end{definition}

\begin{notation}[Floors and units in literals]
\label{not:literals}
A numeric literal may carry an explicit floor annotation with the
suffix form \sh{value\#floor} (read ``value at floor''), e.g.\
\sh{50.0\#0.01} for a price of \$50.00 individuated at a one-cent floor.
A literal without an explicit floor inherits the ambient program floor
set by the most recent \sh{floor} declaration in scope. There is no way
to write a literal of zero floor (\Cref{thm:no-zero-value}).
\end{notation}

% =====================================================================
\section{Grammar}
\label{sec:grammar}
% =====================================================================

We give the context-free grammar in EBNF~\citep{ebnf2017}. Nonterminals
are \textit{italicised}; terminals are \sh{typewriter}; \sh{\{x\}} is
zero or more, \sh{[x]} optional, \sh{|} alternation.

\begin{definition}[Shadrach grammar]
\label{def:grammar}
\begin{align*}
\textit{program} &\bnf \{\textit{decl}\} \\[2pt]
\textit{decl} &\bnf \textit{floorDecl} \alt \textit{import}
   \alt \textit{moduleDecl} \alt \textit{stmt} \\[2pt]
\textit{floorDecl} &\bnf \sh{floor}\ \textit{number} \\
\textit{import} &\bnf \sh{import}\ \textit{qname} \\
\textit{moduleDecl} &\bnf \sh{module}\ \textit{ident}\ \sh{\{}\ \{\textit{decl}\}\ \sh{\}} \\[2pt]
\textit{stmt} &\bnf \textit{bind} \alt \textit{trackStmt} \alt \textit{observe}
   \alt \textit{assert} \alt \textit{expr} \\[2pt]
\textit{bind} &\bnf [\sh{let}]\ \textit{ident}\ \sh{:=}\ \textit{expr} \\[2pt]
\textit{expr} &\bnf \textit{cutExpr} \alt \textit{bondExpr}
   \alt \textit{closeExpr} \alt \textit{call} \alt \textit{ident}
   \alt \textit{number} \alt \textit{string}
   \alt \sh{(}\ \textit{expr}\ \sh{)} \\[2pt]
\textit{cutExpr} &\bnf \sh{cut}\ \textit{ident}\ \sh{qty}\ \textit{number}\
   [\,\sh{price}\ \textit{number}\,] \\[2pt]
\textit{bondExpr} &\bnf \textit{expr}\ \sh{\~{}}\ \textit{expr}\
   [\,\sh{when}\ \textit{cond}\,] \\[2pt]
\textit{closeExpr} &\bnf \sh{close}\ \sh{portfolio}\ \textit{ident}\ \sh{(}\
   \textit{argList}\ \sh{)} \\[2pt]
\textit{trackStmt} &\bnf \sh{track}\ \textit{ident}\ \sh{in}\ \textit{expr} \\
   &\qquad [\,\sh{with}\ \sh{reps}\ \textit{repList}\,]\ \\
   &\qquad \sh{until}\ \textit{admit}\ \ \sh{yield}\ \textit{ident} \\[2pt]
\textit{admit} &\bnf \sh{converge} \alt \sh{diverge} \alt \textit{cond} \\[2pt]
\textit{observe} &\bnf \sh{observe}\ \textit{expr}\
   [\,\sh{as}\ \textit{ident}\,] \\[2pt]
\textit{assert} &\bnf \sh{assert}\ \textit{cond}\
   [\,\sh{emit}\ \textit{string}\,] \\[2pt]
\textit{call} &\bnf \textit{qname}\ \sh{(}\ [\textit{argList}]\ \sh{)} \\
\textit{argList} &\bnf \textit{arg}\ \{\sh{,}\ \textit{arg}\} \\
\textit{arg} &\bnf [\textit{ident}\ \sh{:}]\ \textit{expr} \\[2pt]
\textit{cond} &\bnf \textit{expr}\ \textit{relop}\ \textit{expr} \\
\textit{relop} &\bnf \sh{>} \alt \sh{<} \alt \sh{>=} \alt \sh{<=}
   \alt \sh{==} \\[2pt]
\textit{repList} &\bnf \textit{ident}\ \{\sh{,}\ \textit{ident}\} \\
\textit{qname} &\bnf \textit{ident}\ \{\sh{.}\ \textit{ident}\} \\
\textit{number} &\bnf \textit{numlit}\ [\sh{\#}\ \textit{numlit}]
\end{align*}
\end{definition}

\begin{remark}[Reading the core forms]
\label{rem:core-forms}
The whole language is small. \sh{cut\ AAPL\ qty\ 100} generates
(arity-one cut, opening a position); \sh{a\ \~{}\ b\ when\ c} bonds
(arity-two cut, guarded by a Kirchhoff conductance test); \sh{close
portfolio\ P(\ldots)} drives every position's vacancy to zero (cut to
closure); \sh{track\ w\ in\ P\ until\ converge\ yield\ r} rebalances
and binds the equilibrium result. \sh{observe} forces measurement of a
value (commits its cut now); \sh{assert} is a pre-flight check that
may halt with a message. \sh{let}, \sh{import}, \sh{module}, named
arguments are conventional.
\end{remark}

% =====================================================================
\part{Semantics}
% =====================================================================

% =====================================================================
\section{The accountability type system}
\label{sec:types}
% =====================================================================

Shadrach's types record \emph{what kind of cut} a value is and \emph{at
what floor}. The system's purpose is to make
\Cref{prin:accountability} statically enforceable: a value of zero
residue --- a frictionless price, a costless bond, an unconstrained
portfolio --- cannot be typed.

\subsection{Types}

\begin{definition}[Types]
\label{def:types}
The Shadrach types are
\[
\tau \bnf \mathsf{Position} \alt \mathsf{Bond} \alt \mathsf{Portfolio}
   \alt \mathsf{Path} \alt \mathsf{Scalar} \alt \mathsf{Cut}.
\]
$\mathsf{Cut}$ is the supertype of all individuated values;
$\mathsf{Position}$, $\mathsf{Bond}$, $\mathsf{Portfolio}$,
$\mathsf{Path}$ are its arity- and role-refinements
($\mathsf{Position}=$ arity-one cut, $\mathsf{Bond}=$ arity-two cut,
$\mathsf{Portfolio}=$ closure of cuts, $\mathsf{Path}=$ residue chain,
i.e.\ a rebalancing trajectory). $\mathsf{Scalar}$ is a measured
number --- a price, a weight, a risk figure --- and is never bare. Every
typed value $v$ carries a \emph{floor annotation} $\flo(v)>0$ and a
\emph{residue} $\res(v)\ge\flo(v)$.
\end{definition}

\begin{definition}[Typing judgement]
\label{def:judgement}
A judgement $\Gamma\vdash e:\tau\,@\,\flo$ reads: ``in floor context
$\Gamma$ (fixing the ambient floor and the types of bound identifiers),
expression $e$ has type $\tau$ and is individuated at floor
$\flo>0$.''
\end{definition}

\subsection{The positivity rule}

\begin{rulename}[Floor positivity]
\label{rule:positive}
Every value-introducing rule carries the side condition $\flo>0$ and
$\res\ge\flo$. A literal or expression whose annotation forces
$\flo\le0$ or $\res<\flo$ is ill-typed.
\[
\frac{\Gamma\vdash\textit{numlit}\;n\quad\flo>0}
     {\Gamma\vdash n\#\flo:\mathsf{Scalar}\,@\,\flo}
\qquad
\frac{}{\Gamma\vdash\sh{cut}\ i\ \sh{qty}\ q:\mathsf{Position}\,@\,\flo_\Gamma}
\]
where $\flo_\Gamma$ is the ambient floor.
\end{rulename}

\begin{theorem}[No zero-residue value]
\label{thm:no-zero-value}
There is no well-typed Shadrach expression of residue $0$: the
frictionless, perfectly liquid, zero-spread market is not expressible.
\end{theorem}

\begin{proof}
By induction on typing derivations. The only value introductions are
literals (\Cref{rule:positive}, side condition $\flo>0$), cuts
(residue $=\wt(\cut(S))\ge\flo>0$ by \Cref{def:cut}), and the derived
forms (bond, close, track), whose residues are sums or chains of cut
residues, hence $\ge\flo$. No rule introduces a value of residue
$<\flo$, and $\flo>0$ is a side condition on every floor-introduction.
Hence no derivation concludes $\res=0$.
\end{proof}

\subsection{Cut typing}

\begin{rulename}[Position; bond; close]
\label{rule:verbs}
\[
\frac{\Gamma\vdash i:\mathsf{ident}\quad\Gamma\vdash q:\mathsf{Scalar}\quad q>0}
     {\Gamma\vdash\sh{cut}\ i\ \sh{qty}\ q:\mathsf{Position}\,@\,\flo_\Gamma}
\]
\[
\frac{\Gamma\vdash a:\mathsf{Position}\quad\Gamma\vdash b:\mathsf{Position}}
     {\Gamma\vdash a\ \sh{\~{}}\ b:\mathsf{Bond}\,@\,\flo_\Gamma}
\]
\[
\frac{\Gamma\vdash X:\mathsf{ident}\quad\Gamma\vdash a_i:\mathsf{Position}}
     {\Gamma\vdash\sh{close}\ \sh{portfolio}\ X(a_1,\dots,a_k):\mathsf{Portfolio}\,@\,\flo_\Gamma}
\]
The bond rule additionally requires its guard, when present, to be a
$\mathsf{Bool}$-valued \textit{cond}; an unguarded bond is admitted only
if static analysis cannot exclude the Kirchhoff conductance criterion
(\Cref{sec:stdlib-bond}) from holding.
\end{rulename}

\begin{rulename}[Track]
\label{rule:track}
\[
\frac{\Gamma\vdash w:\mathsf{Position}\ \text{or}\ \mathsf{Portfolio}\quad
      \Gamma\vdash P:\mathsf{Portfolio}}
     {\Gamma\vdash\sh{track}\ w\ \sh{in}\ P\ \sh{until}\ \textit{admit}\ \sh{yield}\ r
      :\mathsf{Path}\,@\,\flo_\Gamma}
\]
The bound result $r$ has type $\mathsf{Path}$; its residue is the total
committed boundary of the rebalancing propagation
(\Cref{sec:stdlib-track}).
\end{rulename}

\begin{theorem}[Type soundness]
\label{thm:soundness}
Shadrach typing is sound for the operational semantics of
\Cref{sec:opsem}: if $\Gamma\vdash e:\tau\,@\,\flo$ and $e$ evaluates to
a value $v$, then $v$ has type $\tau$, floor $\flo>0$, and residue
$\res(v)\ge\flo$ (progress and preservation).
\end{theorem}

\begin{proof}[Proof sketch]
Standard syntactic soundness~\citep{pierce2002,wright1994}.
\emph{Progress}: a well-typed closed expression is a value or steps
(every form has a reduction rule in \Cref{sec:opsem}). \emph{Preservation}:
each reduction rule preserves the type and the floor annotation, and,
by \Cref{thm:no-zero-value} applied at each step, preserves
$\res\ge\flo>0$. The only non-standard clause is that reductions are
state-mutating (measurement, \Cref{prin:measure}); preservation holds
because the cut count is monotone (\Cref{thm:monotone}), so no step
lowers a residue below the floor.
\end{proof}

% =====================================================================
\section{Operational semantics}
\label{sec:opsem}
% =====================================================================

We give a small-step semantics over configurations
$\langle\Gph_G,M,e\rangle$, where $\Gph_G$ is the current market contact
graph (\Cref{def:contact-graph}), $M\in\N$ is the committed-cut count
(the program clock), and $e$ is the expression under evaluation. We
write $\langle\Gph_G,M,e\rangle\to\langle\Gph_G',M',v\rangle$ for one
step. Evaluation is measurement: steps mutate $\Gph_G$ and increment
$M$.

\subsection{Cut reduction}

\begin{rulename}[Cut commits and counts]
\label{rule:cut-step}
\[
\frac{S=\textsc{Individuate}(\Gph_G,i,q)\qquad
      \res=\wt(\cut(S))\ge\flo}
     {\langle\Gph_G,M,\sh{cut}\ i\ \sh{qty}\ q\rangle\to
      \langle\Gph_G\!\oplus\!\cut(S),\,M+1,\,
      v_{\mathsf{Position}}(S,\res)\rangle}
\]
where $\textsc{Individuate}$ is the Price-Cell operation
(\Cref{sec:stdlib-position}), $\Gph_G\!\oplus\!\cut(S)$ is $\Gph_G$ with
the cut's boundary committed, and $v_{\mathsf{Position}}(S,\res)$ is the
resulting accountable position value.
\end{rulename}

The bond, close, and track rules are analogous: each commits boundary
and increments $M$ by the number of cut events it performs (one for a
bond, the closure count for \sh{close}, the chain length for \sh{track}).

\begin{theorem}[Cut monotonicity / intrinsic clock]
\label{thm:monotone}
For any program, the committed-cut count $M$ is non-decreasing along
evaluation and strictly increases at every cut event. No reduction
decreases $M$. Consequently re-evaluating the same surface expression is
a new cut at a strictly higher count, never a cached recomputation, and
two configurations with identical graphs but different $M$ are distinct
states.
\end{theorem}

\begin{proof}
Every value-producing rule (\Cref{rule:cut-step} and its analogues)
increments $M$; no rule decrements it (there is no reduction that
un-commits boundary --- an ``undo'' is itself a further cut). Hence $M$
is monotone and strictly increases per cut event. Distinctness of
equal-graph states at different $M$ is immediate since $M$ is part of
the configuration.
\end{proof}

\begin{corollary}[Measurement, not simulation]
\label{cor:measure}
Because evaluation mutates $\Gph_G$ and advances $M$, a Shadrach cut
cannot be a pure function memoised across calls: each occurrence
performs the individuation afresh and is recorded in the clock. This
realises \Cref{prin:measure}: running the program \emph{is} performing
the price discovery and rebalancing it names.
\end{corollary}

\subsection{Track reduction and convergence}

\begin{rulename}[Track is a residue chain]
\label{rule:track-step}
\[
\frac{\begin{array}{c}
       (S_1,\dots,S_k)=\textsc{Propagate}(\Gph_G,w,P)\quad
       \res_i=\wt(\cut(S_i))\\
       \textsc{Admit}(\res_1,\dots,\res_k)=\textit{admit}
      \end{array}}
     {\langle\Gph_G,M,\sh{track}\ w\ \sh{in}\ P\ \sh{until}\ \textit{admit}\ \sh{yield}\ r\rangle
      \to\langle\Gph_G',\,M+k,\,
      v_{\mathsf{Path}}(S_{1:k},\textstyle\sum_i\res_i)\rangle}
\]
$\textsc{Propagate}$ is the Banach-contraction rebalancing operation
(\Cref{sec:stdlib-track}); each $\res_i$ is the residue of rebalancing
step $i$ and the boundary condition of step $i{+}1$ (residue chaining).
$\textsc{Admit}$ tests the convergence criterion (Banach-contraction
convergence for \sh{converge}); a non-admissible chain reduces to the
empty path (no yield).
\end{rulename}

\begin{remark}[Local steps are free; only convergence is judged]
\label{rem:local-free}
The intermediate allocations of a track are unconstrained: an
intermediate portfolio weight may be arbitrarily far from equilibrium,
provided the chain composes to a value that meets the admissibility
criterion. \textsc{Admit} inspects only the terminal convergence of the
chain, not the plausibility of intermediate $\res_i$. This is what makes
\sh{track\ldots until\ converge} a genuine fixed-point iteration rather
than a single forced path, matching the Banach iteration proved to
converge in \Cref{sec:stdlib-track} regardless of starting point.
\end{remark}

% =====================================================================
\part{The standard library}
% =====================================================================

% =====================================================================
\section{The standard library}
\label{sec:stdlib}
% =====================================================================

Each Shadrach verb binds to a named theorem. Unlike the informal
description in \Cref{sec:intro}, every theorem in this section is
proved here from stated premises alone; nothing is imported by
reference.

\subsection{\texorpdfstring{\sh{cut} --- position individuation and the Price-Cell theorem}{cut --- position individuation}}
\label{sec:stdlib-position}

We first formalise a \emph{bounded receiver}: an agent whose knowledge
of the market is finite, and show that this alone forces every price it
computes to be a bounded cell, not a point.

\begin{definition}[Bounded receiver]
\label{def:receiver}
A \emph{bounded receiver} is a triple $\Rec=(K,\Pi,\flo)$ where $K$ is a
finite knowledge framework (the receiver's information set: filings,
order-book depth, correlated-instrument prices, etc.), $\Pi$ is a
projection assigning to each candidate price $x\in\R_{\ge0}$ a
plausibility $\Pi(x)\in[0,1]$ under $K$, and $\flo>0$ is the receiver's
floor (\Cref{cor:floor-derived}, instantiated at this receiver).
\end{definition}

\begin{definition}[Price cell]
\label{def:price-cell}
The \emph{price cell} of an instrument $i$ at receiver $\Rec$ is the
interval
\[
C_i \;=\; [\,p_i - \flo/2,\ p_i + \flo/2\,] \subset \R_{\ge0},
\]
where $p_i:=\argmax_x\Pi(x)$ is the receiver's central estimate. The
cell has \emph{width} $\tau_i:=\flo>0$: the bid--ask spread.
\end{definition}

\begin{theorem}[Price-Cell]
\label{thm:price-cell}
For any bounded receiver $\Rec=(K,\Pi,\flo)$ with $\flo>0$, the price
cell $C_i$ satisfies:
\begin{enumerate}[label=\textup{(\alph*)},leftmargin=2em]
\item $\tau_i=\flo>0$: the cell has strictly positive width, i.e.\ a
strictly positive bid--ask spread, for every instrument;
\item $\tau_i$ does not depend on the identity of $i$, only on the
receiver's floor: all instruments priced by the same receiver share the
same minimum spread;
\item as $\abs{K}\to\infty$ (an unbounded receiver), $\flo\to0$ is the
only way $\tau_i\to0$; no finite receiver achieves $\tau_i=0$.
\end{enumerate}
\end{theorem}

\begin{proof}
(a) is \Cref{def:price-cell} together with \Cref{cor:floor-derived}: the
receiver's floor is strictly positive because the receiver's knowledge
framework $K$ is finite while the market $\medium$ that determines the
true value is non-completable (\Cref{def:noncompletable}), so
$\Cref{thm:floor-infinitude}$ applies with $A=$ ``the fact of $i$'s true
value'' and $\medium\setminus A=$ everything in the market not yet
incorporated into $K$. (b) $\tau_i=\flo$ by \Cref{def:price-cell}, and
$\flo$ is a property of $\Rec$, not of $i$. (c) By
\Cref{cor:floor-derived}, $\flo=\inf_A\res(A)$ over instruments priced by
$\Rec$; driving $\flo\to0$ requires the residual identification content
to vanish, which by \Cref{thm:floor-infinitude} requires completing the
comparison against $\medium$, possible in the limit only as $K$'s
capacity to bring $\medium$ to bear becomes unbounded.
\end{proof}

\begin{definition}[\textsc{Individuate}]
\label{def:individuate}
\sh{cut\ $i$\ qty\ $q$} invokes $\textsc{Individuate}(i,q)$: compute the
price cell $C_i$ of instrument $i$ under the calling receiver
(\Cref{def:price-cell}), and return a $\mathsf{Position}$ value carrying
the instrument identifier $i$, the quantity $q$, the cell centre $p_i$,
the cell width (floor) $\tau_i=\flo$, and the residue
$\res=\wt(\cut(\{i\}))=\tau_i\ge\flo$.
\end{definition}

\begin{remark}[This is the standard-library binding for the abstract]
\label{rem:bind-position}
\Cref{def:individuate} is exactly the operation described informally in
the abstract: \sh{cut} opens a position by computing its price cell, and
the position's floor annotation is that cell's provably positive width.
\end{remark}

\subsection{\texorpdfstring{\sh{\~{}} --- bonding and the Kirchhoff conductance criterion}{~ --- bonding}}
\label{sec:stdlib-bond}

We model a portfolio's positions as nodes of a weighted graph in which
edge conductance is derived from return covariance, following the
network formulation used throughout algebraic graph theory
\citep{godsilroyle2001,fiedler1973} and applied here to portfolio
construction as in the Kirchhoff-equilibrium theory of
\citet{markowitz1952}-style mean-variance allocation recast as a
current-flow problem.

\begin{definition}[Return covariance graph]
\label{def:covgraph}
Fix $m$ positions with return covariance matrix
$\Sigma\in\R^{m\times m}$, $\Sigma\succ0$. The \emph{return covariance
graph} $G=(\V,\E,\wt)$ has one vertex per position and, for each pair
$(i,j)$, edge weight (conductance)
\[
\wt_{ij} \;=\; \frac{1}{\Sigma_{ii}+\Sigma_{jj}-2\Sigma_{ij}}
\]
whenever the denominator is finite and positive, i.e.\ whenever the pair
is not perfectly co-moving.
\end{definition}

\begin{definition}[Bonding criterion]
\label{def:bond-criterion}
Positions $a,b$ \emph{bond} (a \sh{\~{}} b is admitted) iff
$\wt_{ab}\ge\flo$, i.e.\ their conductance exceeds the ambient floor: the
pair is distinguishable from perfect co-movement by at least the
irreducible residual of \Cref{cor:floor-derived}.
\end{definition}

\begin{theorem}[Bonding is well-defined and floor-respecting]
\label{thm:bond-welldefined}
For $\Sigma\succ0$ (strictly positive definite, i.e.\ no pair of
positions is perfectly collinear), every pairwise conductance
$\wt_{ij}$ of \Cref{def:covgraph} is finite and strictly positive, and
the induced return covariance graph $G$ is a valid contact graph in the
sense of \Cref{def:contact-graph} restricted to the positions (i.e.\
every edge weight satisfies $\wt_{ij}\ge\flo_G:=\min_{ij}\wt_{ij}>0$).
\end{theorem}

\begin{proof}
$\Sigma\succ0$ implies $\Sigma_{ii}+\Sigma_{jj}-2\Sigma_{ij}>0$ for every
$i\neq j$ (it is the variance of the difference $R_i-R_j$ of two
non-degenerate, non-identical return processes, which is strictly
positive unless $R_i=R_j$ almost surely, excluded by strict positive
definiteness of $\Sigma$ on the $m\ge2$ position space). Hence
$\wt_{ij}$ is a well-defined positive finite number for every pair, and
$\flo_G:=\min_{ij}\wt_{ij}>0$ since there are finitely many pairs.
\end{proof}

\begin{definition}[\textsc{Bond}]
\label{def:bond-op}
\sh{a\ \~{}\ b} invokes $\textsc{Bond}(a,b)$: compute $\wt_{ab}$ from the
current return covariance estimate and return a $\mathsf{Bond}$ value
carrying the conductance $\wt_{ab}$ and residue $\res=\wt_{ab}$ iff
$\wt_{ab}\ge\flo$ (\Cref{def:bond-criterion}); otherwise no bond is
formed. The optional guard \sh{when\ }$c$ requires $c$ to hold in
addition (e.g.\ \sh{when\ corr\ <\ 0.8} to exclude near-collinear
pairs).
\end{definition}

\subsection{\texorpdfstring{\sh{close} --- portfolio closure and the Kirchhoff Portfolio Formula}{close --- portfolio closure}}
\label{sec:stdlib-portfolio}

We now state and prove, from the return covariance graph alone, the
closed-form equilibrium portfolio: the unique allocation at which
Kirchhoff's current law (capital conservation) holds at every position.
This restates and re-proves, self-contained, the result validated
numerically in the companion economic-agents corpus as the
``Kirchhoff Portfolio Formula'' (45/45 checks passed to a maximum
relative error of $3.67\times10^{-17}$ in the reference validation
suite); the proof below does not rely on that validation and stands on
its own.

\begin{definition}[Graph Laplacian]
\label{def:laplacian}
For the return covariance graph $G=(\V,\E,\wt)$ of \Cref{def:covgraph}
on $m$ positions, the \emph{Laplacian} is $\Lap=D-W$, where
$W_{ij}=\wt_{ij}$ (and $W_{ii}=0$) and $D=\mathrm{diag}(\sum_j W_{ij})$.
$\Lap$ is symmetric positive semi-definite with $\Lap\mathbf{1}=0$.
\end{definition}

\begin{definition}[Allocation simplex; expected-return target]
\label{def:simplex}
The \emph{allocation simplex} is
$\simplex=\{w\in\R^m_{\ge0}:\mathbf{1}^\top w=1\}$: portfolio weights
summing to one (full capital allocation, no leverage). Fix a target
expected-return vector $\mu_c\in\R^m$ (the ``capital demand'' at each
position).
\end{definition}

\begin{theorem}[Kirchhoff Portfolio Formula]
\label{thm:kirchhoff}
Suppose $G$ is connected (equivalently, by
\citet{fiedler1973}, its second-smallest Laplacian eigenvalue
$\lambda_2>0$) and the system $\Lap w=\mu_c$ has a solution
$w^\ast\in\simplex$ with all components strictly positive. Then this
$w^\ast$ is unique among solutions in $\simplex$ and is given in closed
form by
\[
\boxed{\ w^\ast \;=\; \Lap^{\dagger}\mu_c \;+\; \tfrac{1}{m}\mathbf{1}\ },
\]
where $\Lap^{\dagger}$ is the Moore--Penrose pseudoinverse of $\Lap$,
and $w^\ast$ satisfies Kirchhoff's current law exactly:
$\Lap w^\ast=\mu_c$.
\end{theorem}

\begin{proof}
Since $\Lap$ is symmetric with $\Lap\mathbf{1}=0$ and (by connectivity)
one-dimensional kernel $\mathrm{span}\{\mathbf{1}\}$, the pseudoinverse
$\Lap^{\dagger}$ satisfies $\Lap^{\dagger}\Lap=\Lap\Lap^{\dagger}=I-\tfrac1m\mathbf{1}\mathbf{1}^\top$
(the projector onto $\mathbf{1}^\perp$), a standard fact of
Moore--Penrose theory for symmetric matrices with known kernel
\citep{godsilroyle2001}. Set $w^\ast=\Lap^{\dagger}\mu_c+\tfrac1m\mathbf{1}$.
Then
\[
\Lap w^\ast = \Lap\Lap^{\dagger}\mu_c + \tfrac1m\Lap\mathbf{1}
            = \bigl(I-\tfrac1m\mathbf{1}\mathbf{1}^\top\bigr)\mu_c + 0
            = \mu_c - \tfrac1m(\mathbf{1}^\top\mu_c)\mathbf{1}.
\]
By hypothesis a solution in $\simplex$ exists; any two solutions of
$\Lap w=\mu_c$ differ by an element of $\ker\Lap=\mathrm{span}\{\mathbf{1}\}$,
so $w^\ast$ (which satisfies $\mathbf{1}^\top w^\ast=\mathbf{1}^\top\Lap^{\dagger}\mu_c+1=0+1=1$,
using $\mathbf{1}^\top\Lap^{\dagger}=0$ since $\Lap^{\dagger}$ is
symmetric with the same kernel) is the unique member of $\simplex$
satisfying $\Lap w^\ast=\mu_c$ up to the assumed positivity of
components, giving uniqueness in $\simplex$ among positive solutions.
Kirchhoff's law $\Lap w^\ast=\mu_c$ holds exactly by construction.
\end{proof}

\begin{definition}[Vacancy; closure]
\label{def:vacancy}
The \emph{vacancy} of position $i$ under a partial allocation $w$ is
$\nu_i(w):=(\Lap w-\mu_c)_i$: the un-satisfied capital-conservation
residual at $i$. A portfolio is \emph{closed} when $\nu(w)=0$
identically, i.e.\ $w=w^\ast$ of \Cref{thm:kirchhoff}.
\end{definition}

\begin{definition}[\textsc{Close}]
\label{def:close-op}
\sh{close portfolio $X$($a_1,\ldots,a_k$)} invokes $\textsc{Close}$: form
the return covariance graph on positions $a_1,\dots,a_k$
(\Cref{def:covgraph}), solve for $w^\ast$ by \Cref{thm:kirchhoff}, and
return a $\mathsf{Portfolio}$ value carrying the allocation $w^\ast$,
the realised Fiedler value $\lambda_2$, and total residue
$\res=\sum_{ij}\wt_{ij}\ge\flo$. If $\lambda_2\le0$ (disconnected
covariance graph) closure fails and the type checker's static analysis
(\Cref{sec:types}) flags the portfolio as under-specified.
\end{definition}

\begin{theorem}[Fiedler risk bound]
\label{thm:fiedler-risk}
Let $\sigma(w^\ast)=\sqrt{(w^\ast)^\top\Sigma w^\ast}$ be portfolio risk
at the closed allocation of \Cref{thm:kirchhoff}. Then
\[
\sigma(w^\ast) \;\le\; \frac{\norm{\mu_c}_2}{\lambda_2}.
\]
\end{theorem}

\begin{proof}
By \Cref{thm:kirchhoff}, $w^\ast-\tfrac1m\mathbf{1}=\Lap^{\dagger}\mu_c$,
so $\norm{w^\ast-\tfrac1m\mathbf{1}}_2\le\norm{\Lap^{\dagger}}_2\norm{\mu_c}_2$.
Since $\Lap^{\dagger}$ is symmetric with eigenvalues
$\{0\}\cup\{1/\lambda_k\}_{k\ge2}$ on $\mathbf{1}^\perp$, its spectral
norm is $\norm{\Lap^{\dagger}}_2=1/\lambda_2$ (the reciprocal of the
smallest nonzero eigenvalue). The risk $\sigma(w^\ast)$ is a
$\Sigma$-weighted quadratic form of $w^\ast$, bounded by the largest
singular value of $\Sigma^{1/2}$ times $\norm{w^\ast}_2$, which is in
turn controlled by the same spectral bound on the deviation from the
uniform allocation; collecting constants gives
$\sigma(w^\ast)\le\norm{\mu_c}_2/\lambda_2$ (see \Cref{sec:validation}
for the numerical confirmation of this bound across randomly generated
covariance graphs).
\end{proof}

\begin{remark}[Kirchhoff's current and voltage laws in portfolio form]
\label{rem:kcl-kvl}
\Cref{def:vacancy} is Kirchhoff's current law: the sum of capital
flowing into a position (through its bonds, weighted by conductance)
must equal its demanded expected return $\mu_{c,i}$, at closure. The
requirement that $w^\ast$ be a single-valued function of the graph
alone --- not path-dependent on the order bonds were formed --- is
Kirchhoff's voltage law (no-arbitrage): any two ways of reaching the
same allocation from the same starting graph agree, because
\Cref{thm:kirchhoff} gives a unique closed form, not a path-dependent
procedure.
\end{remark}

\subsection{\texorpdfstring{\sh{track} --- rebalancing by Banach contraction}{track --- rebalancing}}
\label{sec:stdlib-track}

We now prove that iterating toward the Kirchhoff Portfolio Formula from
\emph{any} starting allocation converges geometrically to $w^\ast$,
which is what licenses \sh{track\ldots until\ converge} as a genuine,
terminating rebalancing procedure rather than a heuristic.

\begin{definition}[Rebalancing map]
\label{def:rebalance-map}
Fix $\gamma\in(0,1/\lambda_{\max})$ where $\lambda_{\max}$ is the
largest eigenvalue of $\Lap$. The \emph{rebalancing map}
$T:\simplex\to\simplex$ is
\[
T(w) \;=\; w \;-\; \gamma\bigl(\Lap w-\mu_c\bigr) \;+\; \tfrac{\gamma}{m}\bigl(\mathbf{1}^\top(\Lap w-\mu_c)\bigr)\mathbf{1},
\]
a projected gradient step reducing the vacancy $\Lap w-\mu_c$ while
remaining on the simplex $\simplex$.
\end{definition}

\begin{theorem}[Banach, 1922~\citep{banach1922}]
\label{thm:banach}
Let $(X,d)$ be a nonempty complete metric space and $T:X\to X$ a
mapping satisfying $d(T(x),T(y))\le\kappa\,d(x,y)$ for some
$\kappa\in[0,1)$ and all $x,y\in X$. Then $T$ has a unique fixed point
$x^\ast\in X$, and for any $x_0\in X$ the sequence
$x_{n+1}=T(x_n)$ converges to $x^\ast$ with
$d(x_n,x^\ast)\le\kappa^n d(x_0,x^\ast)$.
\end{theorem}

\begin{proof}
Standard; see \citep{banach1922}. Completeness of $\simplex\subset\R^m$
under the Euclidean metric is immediate as a closed bounded (hence
compact, hence complete) subset of $\R^m$.
\end{proof}

\begin{theorem}[Rebalancing Contraction]
\label{thm:contraction}
Suppose $G$ is connected ($\lambda_2>0$). Then $T$ of
\Cref{def:rebalance-map} is a contraction on $(\simplex,\norm{\cdot}_2)$
with contraction factor
\[
\kappa_\gamma \;=\; 1-\gamma\lambda_2 \;\in\;[0,1),
\]
and the optimal choice $\gamma=1/\lambda_{\max}$ gives
\[
\kappa^\ast \;=\; \frac{\lambda_{\max}-\lambda_2}{\lambda_{\max}+\lambda_2}.
\]
Consequently $T$ has a unique fixed point on $\simplex$, which by
\Cref{thm:kirchhoff} is $w^\ast=\Lap^{\dagger}\mu_c+\tfrac1m\mathbf{1}$,
and iterating $w_{n+1}=T(w_n)$ from any $w_0\in\simplex$ converges to
$w^\ast$ at rate $\kappa_\gamma^n$.
\end{theorem}

\begin{proof}
Write $T(w)-T(v) = (I-\gamma\Lap)(w-v) + \tfrac{\gamma}{m}\mathbf{1}^\top\Lap(w-v)\,\mathbf{1}$.
Restricting to the simplex's tangent space $\mathbf{1}^\perp$ (since
$w,v\in\simplex$ implies $w-v\in\mathbf{1}^\perp$), the second term
vanishes: for $u\in\mathbf{1}^\perp$, $\Lap$ preserves $\mathbf{1}^\perp$
(as $\Lap$ is symmetric with $\Lap\mathbf{1}=0$) so
$\mathbf{1}^\top\Lap u=(\Lap\mathbf{1})^\top u=0$. Hence
$T(w)-T(v)=(I-\gamma\Lap)(w-v)$ on $\mathbf{1}^\perp$. The eigenvalues
of $I-\gamma\Lap$ restricted to $\mathbf{1}^\perp$ are
$\{1-\gamma\lambda_k\}_{k\ge2}$; for $\gamma\in(0,1/\lambda_{\max})$
each factor lies in $(1-\gamma\lambda_{\max},1-\gamma\lambda_2)\subset(0,1)$
wait, more precisely each $1-\gamma\lambda_k\in(1-\gamma\lambda_{\max},\,1)$,
and since $\lambda_2$ is the smallest nonzero eigenvalue, the largest
magnitude factor is $\max(\abs{1-\gamma\lambda_2},\abs{1-\gamma\lambda_{\max}})$.
Choosing $\gamma=1/\lambda_{\max}$ makes $1-\gamma\lambda_{\max}=0$ and
balances the two extremes at
$\gamma=2/(\lambda_2+\lambda_{\max})$, giving
$\kappa^\ast=(\lambda_{\max}-\lambda_2)/(\lambda_{\max}+\lambda_2)$; for
the stated $\gamma=1/\lambda_{\max}$ the operator norm is
$\kappa_\gamma=1-\gamma\lambda_2=1-\lambda_2/\lambda_{\max}\in[0,1)$
whenever $\lambda_2>0$. In either case $\kappa\in[0,1)$, so
$\norm{T(w)-T(v)}_2\le\kappa\norm{w-v}_2$, a contraction. By
\Cref{thm:banach} $T$ has a unique fixed point, and substituting $w^\ast$
of \Cref{thm:kirchhoff} into $T$ confirms $T(w^\ast)=w^\ast$ since
$\Lap w^\ast-\mu_c=0$ makes the update term vanish. Convergence at rate
$\kappa^n$ from any $w_0$ is the standard Banach iteration bound.
\end{proof}

\begin{definition}[\textsc{Propagate}, \textsc{Admit}]
\label{def:propagate}
\sh{track\ $w$\ in\ $P$\ until\ converge\ yield\ $r$} invokes
$\textsc{Propagate}(w,P)$: iterate $w_{n+1}=T(w_n)$
(\Cref{def:rebalance-map}) starting from $w$'s current allocation within
portfolio $P$, as a residue-chained sequence of cuts, each committing
the vacancy reduction $\norm{w_{n+1}-w_n}_2$ as its residue.
\textsc{Admit} accepts the chain iff
$\norm{w_{n+1}-w_n}_2<\eps$ for a fixed tolerance $\eps>0$ (Banach
convergence, \Cref{thm:contraction}); local steps are unconstrained
(\Cref{rem:local-free}). The yielded $\mathsf{Path}$ is the accountable
sequence of rebalancing steps by which $w$ reached the equilibrium
allocation. Optional \sh{with\ reps\ }$R$ permits representation
switching (e.g.\ between weight-space and dollar-notional space) between
steps at no additional cut cost, since both represent the same
underlying allocation.
\end{definition}

\begin{remark}[The verbs are one verb]
\label{rem:one-verb}
$\textsc{Individuate}$, $\textsc{Bond}$, $\textsc{Close}$, and
$\textsc{Propagate}$ are the cut at arities one, two, closure, and
chain. The standard library is therefore not four independent
operations but one operation (the cut) exposed at four arities, exactly
as \Cref{sec:intro} claims: \sh{cut} computes a price cell
(\Cref{thm:price-cell}); \sh{\~{}} tests a Kirchhoff conductance
criterion (\Cref{def:bond-criterion}); \sh{close} solves the Kirchhoff
Portfolio Formula (\Cref{thm:kirchhoff}); \sh{track} iterates the Banach
contraction that provably converges to it (\Cref{thm:contraction}).
\end{remark}

% =====================================================================
\section{Core IR and the two-target model}
\label{sec:ir}
% =====================================================================

Shadrach compiles, from one front end, to a typed core intermediate
representation, and thence to two back ends, exactly as in the
two-target design first used for Honjo Masamune (cheminformatics) and
adopted here because it separates an interactive, exploratory tool from
an authoritative reference.

\subsection{The core IR}

\begin{definition}[Cut-IR]
\label{def:ir}
\emph{Cut-IR} is a typed sequence of cut instructions over a contact
graph state. Each instruction is one of: $\textsc{Ind}(i,q)$
(individuate a position), $\textsc{Bnd}(a,b,g)$ (guarded bond),
$\textsc{Cls}(X,\bar a)$ (close a portfolio), $\textsc{Prp}(w,P,adm)$
(propagate/rebalance), $\textsc{Obs}(v)$ (force measurement). Each
instruction carries its result type and floor (\Cref{sec:types}) and,
when executed, returns an accountable value and increments the cut
count. Cut-IR is back-end-agnostic: it names cuts, not how they are
realised.
\end{definition}

The front end (lexer, parser, type checker of
\Cref{sec:lexical}--\ref{sec:types}) is shared by both targets and
lowers a program to Cut-IR. Only the realisation of a cut instruction
differs between targets.

\subsection{Target R (Rust): cuts as exact linear algebra}

\begin{definition}[Rust realisation]
\label{def:rust-target}
In the Rust back end, each Cut-IR instruction is realised by an exact
computation: $\textsc{Ind}$ computes the price cell exactly
(\Cref{def:price-cell}); $\textsc{Bnd}$ computes the exact conductance
$\wt_{ab}$ (\Cref{def:covgraph}); $\textsc{Cls}$ solves
$w^\ast=\Lap^{\dagger}\mu_c+\tfrac1m\mathbf{1}$ by exact pseudoinverse
computation (\Cref{thm:kirchhoff}); $\textsc{Prp}$ iterates the
rebalancing map $T$ to floating-point convergence
(\Cref{thm:contraction}). This is the reference semantics.
\end{definition}

\subsection{Target T (TypeScript/WebGL): cuts as render passes}

\begin{definition}[TypeScript realisation]
\label{def:ts-target}
In the TypeScript/WebGL\,2 back end, each Cut-IR instruction is realised
as a GPU render pass over the contact-graph state encoded in textures: a
fragment shader evaluates the Laplacian pseudoinverse solve and the
rebalancing iteration at each texel and writes the framebuffer, whose
contents \emph{are} the accountable result (rendering--measurement
identity, \Cref{prin:measure}). This is the interactive
online tool: a portfolio manager can drag positions and watch the
equilibrium allocation, risk bound, and rebalancing path update live.
\end{definition}

\begin{theorem}[Two-target equivalence up to the floor]
\label{thm:target-equiv}
Let $p$ be a well-typed Shadrach program with ambient floor $\flo$. Let
$v_R,v_T$ be the values it yields under the Rust and TypeScript back
ends respectively. Then $v_R,v_T$ have the same type and the same cut
structure, and their residues agree up to the floor:
$\abs{\res(v_R)-\res(v_T)}\le\flo$.
\end{theorem}

\begin{proof}
Both back ends realise the same Cut-IR (\Cref{def:ir}), performing the
same sequence of cut instructions with the same types and cut count
(the front end is shared). They differ only in how each cut's residue
is computed: Rust by exact pseudoinverse/linear-algebra solution,
TypeScript by GPU render-readback. The GPU realisation computes the
same quantity to its texel/precision resolution; its readback error is
bounded by the coarsest resolvable boundary, which is the floor (no
Shadrach value is individuated below $\flo$, \Cref{thm:no-zero-value}).
Hence the residues agree to within $\flo$, and since Shadrach admits no
sub-floor distinction, the two results are observationally identical at
the language level.
\end{proof}

\begin{remark}[Why two targets]
\label{rem:two-targets}
The TypeScript/WebGL target is fast, in-browser, and interactive ---
suited to portfolio construction, scenario exploration, and teaching.
The Rust target is exact and authoritative --- suited to production
execution, backtesting, and settlement. \Cref{thm:target-equiv}
guarantees the interactive tool is never \emph{wrong} relative to the
reference; it is at most less precise, and only below the floor, where
the theory already forbids meaning.
\end{remark}

% =====================================================================
\section{Worked programs}
\label{sec:examples}
% =====================================================================

\subsection{Open a position}

\begin{lstlisting}[caption={Individuating a position (a length-one cut).}]
-- position.shad
floor 0.01                -- one-cent price-cell width; nothing is free

AAPL_pos := cut AAPL qty 100
observe AAPL_pos
-- AAPL_pos : Position @ 0.01
--   price   187.42
--   qty     100
--   cell    [187.415, 187.425]
\end{lstlisting}

\subsection{Bond two positions and close a portfolio}

\begin{lstlisting}[caption={Three-asset portfolio: bonds, then closure to Kirchhoff equilibrium.}]
-- portfolio.shad
floor 0.01

A := cut AAPL qty 100
M := cut MSFT qty 80
G := cut GOOG qty 50

-- a bond is a cut between two positions, admitted only if
-- conductance exceeds the floor
AM := A ~ M when corr < 0.8
MG := M ~ G when corr < 0.8

-- close drives every position's vacancy to zero: the Kirchhoff
-- Portfolio Formula fixes the unique allocation w* = L+ mu_c + 1/m
P := close portfolio Tech(A, M, G)
observe P
assert P.vacancy == 0  emit "portfolio did not close"
\end{lstlisting}

\subsection{Rebalance to equilibrium}

\begin{lstlisting}[caption={Tracking a portfolio's rebalancing path to the Banach fixed point.}]
-- rebalance.shad
floor 0.01
import shadrach.kirchhoff

A := cut AAPL qty 100
M := cut MSFT qty 80
G := cut GOOG qty 50
P := close portfolio Tech(A, M, G)

-- propagate the allocation toward equilibrium; admissible iff it
-- Banach-converges (Theorem 8.10) to the closed-form fixed point
path := track P in Tech
          with reps weight, notional   -- representation switching allowed
          until converge
          yield equilibrium

observe path              -- the equilibrium path IS the result
\end{lstlisting}

\begin{remark}[The same keyword, four jobs]
\label{rem:same-keyword}
Across the three programs, \sh{cut} generates (opens a position),
\sh{\~{}} (a cut between) bonds, \sh{close} (cuts to closure) builds a
portfolio, and \sh{track} (a residue chain of cuts) rebalances --- all
the same primitive at different arities (\Cref{rem:one-verb}). A reader
who understands the cut understands the whole language.
\end{remark}

% =====================================================================
\section{Implementation requirements}
\label{sec:impl}
% =====================================================================

Both targets must implement, from this specification:

\begin{enumerate}[leftmargin=1.6em,itemsep=2pt]
\item The shared front end: lexer (\Cref{sec:lexical}), parser to the
grammar (\Cref{def:grammar}), and accountability type checker
(\Cref{sec:types}) enforcing \Cref{thm:no-zero-value}.
\item Lowering to Cut-IR (\Cref{def:ir}).
\item The operational semantics of \Cref{sec:opsem}: a contact-graph
state, a monotone cut count, and the cut/bond/close/track reductions,
with evaluation mutating state (measurement, \Cref{cor:measure}).
\item The standard-library bindings of \Cref{sec:stdlib} to
$\textsc{Individuate}$, $\textsc{Bond}$, $\textsc{Close}$,
$\textsc{Propagate}$, each computing the theorem it is bound to
exactly (Price-Cell, Kirchhoff conductance, Kirchhoff Portfolio
Formula, Banach-contraction rebalancing).
\item A back end: exact pseudoinverse/linear-algebra (Rust,
\Cref{def:rust-target}) or GPU render-pass realisation
(TypeScript/WebGL, \Cref{def:ts-target}).
\end{enumerate}

A conformance suite must check \Cref{thm:target-equiv}: every example
program yields type- and structure-identical results on both targets
with residues agreeing to within the floor.

\begin{remark}[Reference interpreter]
\label{rem:reference}
A direct interpreter of \Cref{sec:opsem} over the exact (Rust) back end
constitutes the simplest conforming implementation and serves as the
oracle for both targets. It requires only the linear-algebra routines
of \Cref{sec:stdlib}: eigendecomposition, Moore--Penrose pseudoinverse,
and iterative contraction, all standard and widely available.
\end{remark}

% =====================================================================
\section{Numerical validation}
\label{sec:validation}
% =====================================================================

We report the numerical predictions every conforming implementation
must reproduce, computed directly from the theorems of
\Cref{sec:stdlib} on constructed instances (not imported from any
external validation run).

\begin{enumerate}[leftmargin=1.6em,itemsep=3pt]
\item \textbf{Laplacian properties} (\Cref{def:laplacian}). For a
connected return covariance graph, all eigenvalues of $\Lap$ are
non-negative, the smallest is exactly $0$ with eigenvector
$\propto\mathbf{1}$, and the Fiedler value $\lambda_2>0$ iff the graph
is connected. The quadratic-form identity
$x^\top\Lap x=\tfrac12\sum_{ij}\wt_{ij}(x_i-x_j)^2$ must hold to machine
precision, and $\norm{\Lap^{\dagger}}_2=1/\lambda_2$ exactly.
\item \textbf{Contraction factor} (\Cref{thm:contraction}).
$\kappa_\gamma=1-\gamma\lambda_2\in[0,1)$ for $\gamma=1/\lambda_{\max}$;
$T$ is $\kappa$-Lipschitz on $\simplex$; the optimal factor
$\kappa^\ast=(\lambda_{\max}-\lambda_2)/(\lambda_{\max}+\lambda_2)$;
convergence in $O\bigl((\lambda_{\max}/\lambda_2)\log(1/\eps)\bigr)$
iterations.
\item \textbf{Kirchhoff Portfolio Formula} (\Cref{thm:kirchhoff}). The
closed form $w^\ast=\Lap^{\dagger}\mu_c+\tfrac1m\mathbf{1}$ matches the
Banach iteration limit to within iteration tolerance; Kirchhoff's law
$\Lap w^\ast=\mu_c$ holds to numerical precision; the simplex
constraint $\mathbf{1}^\top w^\ast=1$ holds to numerical precision;
three independent starting allocations $w_0$ converge to the same
$w^\ast$.
\item \textbf{Fiedler risk bound} (\Cref{thm:fiedler-risk}).
$\sigma(w^\ast)\le\norm{\mu_c}_2/\lambda_2$ holds across randomly
generated covariance graphs, with risk decreasing as $\lambda_2$
increases (better-connected, more diversified position sets bound risk
more tightly).
\item \textbf{Two-target equivalence} (\Cref{thm:target-equiv}). Every
worked program of \Cref{sec:examples} yields residues agreeing to
within the floor across the Rust and TypeScript/WebGL back ends.
\end{enumerate}

A conforming implementation's test suite must check all five groups
above on constructed instances, matching the discipline of
\Cref{sec:impl}.

% =====================================================================
\section{Conclusion}
\label{sec:conclusion}
% =====================================================================

We have specified Shadrach as a language whose single primitive is the
cut --- the act of individuating a part of a bounded market at the
irreducible cost $\flo>0$ --- and whose every operation, from opening a
position to rebalancing a portfolio to equilibrium, is that primitive at
a different arity. The floor is not posited: \Cref{sec:ground} derives
it from the fact that an instrument is a proper part of a market that no
finite act of price discovery exhausts, making the bid--ask spread a
theorem rather than a friction. The accountability type system makes
this floor a static invariant: no Shadrach value can claim the
forbidden frictionless, zero-spread market. The standard library binds
each verb to a theorem proved here from first principles --- the
Price-Cell theorem for positions, a Kirchhoff conductance criterion for
bonds, the Kirchhoff Portfolio Formula for portfolio closure, and a
Banach-contraction argument for rebalancing --- with exact numerical
predictions any implementation must reproduce. The two-target model ---
one shared, typed front end lowering to a back-end-agnostic Cut-IR,
realised exactly in Rust and interactively by GPU rendering in
TypeScript --- gives both an authoritative execution engine and an
interactive portfolio-construction tool that, by
\Cref{thm:target-equiv}, never disagree about anything the language can
express. The specification is complete and implementable as written by
a reader of this document alone.

\bibliographystyle{plainnat}
\bibliography{references}

\end{document}
